\answer{ex:07:1}{$V=\kappa p/\bigl(\kappa p+(1-\kappa)(1-p)\bigr)$：$0.059$、$0.20$、$0.50$。中央値は$0$だが、点~\eqref{eq:expectile}は$\kappa$とともになめらかに変わる。}
\answer{ex:07:2}{$p=1/3$、$x=0.75$、$\sigma_r=0.35$、$V(0.2)=0.083$。}
\answer{ex:07:3}{$V=\sqrt\kappa/(\sqrt\kappa+\sqrt{1-\kappa})$で、$0.25$から$0.75$まで。ばらつきは$\propto\sqrt\eta$。2滴では$V=0.25+0.5\kappa$で、両端のニューロンで分布の差は$0.05$なので、$\eta\lesssim0.01$が必要。}
\answer{ex:07:4}{(a)~$J^*=2\sqrt{C\bar R}$。(b)~払いすぎは$\bigl(k^{1/2}+k^{-1/2}\bigr)/2-1$で、$k=2$で$6\,\%$、$k=4$で$25\,\%$。「2倍以内」の精度で足りる。}
\answer{ex:07:5}{面積$R(e^{-1}-e^{-2})\approx0.23R$のバーストで、ランプ$2.3\,\text{s}$分にあたる。ドーパミンが$V$を伝えるならバーストはなく、レベルが階段状に$e$倍になるだけである。}
\answer{ex:07:6}{出来事は重みを$\propto A\tau_f/(\tau_e+\tau_f)$だけ動かし、レベルは誤差$s\tau_f$を生む。$9\,\text{s}\le\tau_f\le17\,\text{s}$。}
\answer{ex:07:7}{初期状態で$\Delta P=0.80$、$M=0.90$。(a)~$\Delta P=0.74$（$-8\,\%$）、$M=0.43$（$-52\,\%$）。(b)~$\Delta P=0.40$（$-50\,\%$）、$M=0.80$（$-11\,\%$）。二重解離である。}
\answer{ex:07:8}{$N_{\text{eff}}\to2+\rho^2N\approx10^6$試行。配線1本の場合の$10^4$分の1だが、$1\,\%$の漏れでもなお100万試行かかる。}
